\documentclass[11pt,a4paper]{article}

%% PACHETE MATEMATICE
\usepackage{amsmath,amssymb,amsfonts}
\usepackage{amsthm}
\usepackage{mathpazo}
\usepackage{eulervm}
\usepackage{enumerate}
\usepackage{color}
\usepackage{graphicx}
\usepackage[all]{xy}
\usepackage{xcolor}
\usepackage{framed}
% \usepackage[unicode]{hyperref} 
\setcounter{MaxMatrixCols}{20}
\usepackage{multicol}
%\usepackage[romanian]{babel}


%% ASPECT
\linespread{1.05}
\usepackage[T1]{fontenc}
\usepackage[utf8x]{inputenc}
\usepackage[margin=2cm]{geometry}
% \usepackage[color]{showkeys}
\usepackage{anyfontsize}
\usepackage{titlesec}
\titleformat*{\section}{\large\bfseries}

\usepackage{fancyhdr}
\pagestyle{fancy}
\rhead{\small \color{gray} Notițe de Adrian Manea}
\lhead{\small \color{gray} Aero-BCD, 2017---2018, M. Olteanu \& L. Costache}


\usepackage[autostyle = false, style = german]{csquotes}
\usepackage[romanian]{babel}
\newcommand\qq{\enquote}

%% COMENZILE MELE
\renewcommand*{\proofname}{Dem.:}
\newcommand\bb{\mathbb}
\newcommand\dr{\mathrm}
\newcommand\kal{\mathcal}
\newcommand\fk{\mathfrak}
\newcommand\op{\oplus}
\newcommand\ot{\otimes}
\newcommand\ds{\displaystyle}
\newcommand\id{\indent}
\newcommand\nid{\noindent}
\newcommand\seq{\subseteq}
\newcommand\sneq{\subsetneq}
\newcommand\speq{\supseteq}
\newcommand\spneq{\supsetneq}
\newcommand\rar{\rightarrow}
\newcommand\Rar{\Rightarrow}
\newcommand\xrar{\xrightarrow}
\newcommand\lar{\leftarrow}
\newcommand\Lar{\Leftarrow}
\newcommand\xlar{\xleftarrow}
\newcommand\lrar{\leftrightarrow}
\newcommand\Lrar{\Leftrightarrow}
\newcommand\ideal{\trianglelefteq}
\newcommand\ai{\text{ a.\^{i}. }}
\newcommand\sk{(\textit{Schi\c{t}\u{a}}) }
\newcommand\rhup{\rightharpoonup}
\newcommand\rhdn{\rightharpoondown}
\newcommand\lhup{\leftharpoonup}
\newcommand\lhdn{\leftharpoondown}
\newcommand\vid{\emptyset}
\newcommand\wh{\widehat}
\newcommand\ol{\overline}
\newcommand\NN{\mathbb{N}}
\newcommand\RR{\mathbb{R}}
\newcommand\ZZ{\mathbb{Z}}
\newcommand\QQ{\mathbb{Q}}
\newcommand\CC{\mathbb{C}}

%% STILURILE MELE
\newtheoremstyle{dotless}{}{}{\itshape}{}{\bfseries}{:}{ }{}

\theoremstyle{dotless}
\newtheorem{theorem}{Teorem\u{a}}[section]
\newtheorem{proposition}{Propozi\c{t}ie}[section]
\newtheorem{lemma}{Lem\u{a}}[section]
\newtheorem{corollary}{Corolar}[section]
\newtheorem{exercise}{Exerci\c{t}iu}

\newtheoremstyle{dotlessdr}{}{}{}{}{\bfseries}{:}{ }{}
\theoremstyle{dotlessdr}
\newtheorem{example}{Exemplu}[section]
\newtheorem{definition}{Defini\c{t}ie}[section]
\newtheorem{remark}{Observa\c{t}ie}[section]




\begin{document}

\begin{center}
  {\large \textbf{Seminar 4 --- Supliment \\ Densitatea lui $\QQ$ în $\RR$}}
\end{center}

\vspace{2cm}

\section{Principiul lui Arhimede}
\label{sec:arhimede}

Vom aborda problema densității lui $\QQ$ în $\RR$ pas cu pas:

\begin{lemma}\label{le:z-nemarg}
  Fie $S \seq \ZZ$ o submulțime nevidă. Presupunem că $S$ este mărginită superior.
  Atunci $S$ are un cel mai mare element.
\end{lemma}

\begin{proof}
  Deoarece $S$ este nevidă și mărginită superior, rezultă că există $w = \sup S$.
  Arătăm că $w \in S$.

  Dacă o mulțime conține supremumul, acela este și cel mai mare element al său, evident. Ne vom
  baza pe această observație pentru a demonstra. Presupunem că $w \notin S$. Deoarece $w - 1 < 2$,
  rezultă că există $m \in S$, cu $w - 1 < m \leq w$. Cum $w \notin S$, inegalitățile pot fi ambele
  luate ca stricte.

  Continuăm, putînd găsi $n \in S$, cu $m < n < w$. Deci:
  \[
    w - 1 < m < n < w.
  \]
  Dar, din $n < w$, avem $-w < - n$, pe care o putem aduna cu cea de mai sus și
  obținem $-1 < m - n$. Deoarece $m < n$, avem $n - m > 0$, deci $0 < n - m < 1$.
  Cum $m, n \in S$, sînt numere întregi, deci $n - m \in \ZZ$, contradicție.

  Concluzia este că $w \in S$.
\end{proof}

\begin{remark}\label{rk:pentru-r}
  De remarcat este faptul că acest argument nu poate funcționa pentru o mulțime
  continuă. În demonstrație, am folosit în mod esențial faptul că orice număr
  întreg are un predecesor și un succesor, lucru care nu este adevărat pentru $\RR$,
  de exemplu.
\end{remark}

Similar se demonstrează și:

\begin{lemma}\label{le:min-elem-z}
  Fie $S \seq \ZZ$ și presupunem $S \neq \emptyset$ și $S$ mărginită inferior.
  Atunci $S$ are un cel mai mic element.
\end{lemma}

Din aceste rezultate, avem:

\begin{theorem}\label{thm:z-nemarg}
  $\ZZ$ este nemărginită și superior, și inferior.
\end{theorem}

\begin{proof}
  Rezultatul reiese imediat din cele două leme de mai sus, deoarece $\ZZ$ nu are
  nici un cel mai mare element, nici un cel mai mic element.
\end{proof}

\begin{theorem}[Principiul lui Arhimede]\label{thm:arhimede}
  Fie $x \in \RR$ și $x > 0$. Atunci există un întreg pozitiv $n$, astfel încît $\frac{1}{n} < x$.
\end{theorem}

\begin{proof}
  Fie $x$ ca în teoremă. Atunci $\frac{1}{x}$ nu poate fi o margine superioară pentru $\ZZ$, deoarece
  $\ZZ$ nu este mărginită superior. Așadar, putem găsi $n \in \ZZ$, chiar $n > 0$, cu $\frac{1}{x} < n$.

  Rezultă $\frac{1}{n} < x$.
\end{proof}

Această teoremă ne arată că, alegînd $n$ suficient de mare, putem face $\frac{1}{n}$
cît de aproape de 0 dorim.

%%%%%%%%%%%%%%%%%%%%%%%%%%%%%%%%%%%%%%%%%%%%%%%%%%%%%%%%%%%%%%%%%%%%%%

\section{Densitatea lui $\QQ$ în $\RR$}
\label{sec:densitate}

\begin{definition}\label{def:densitate}
  O submulțime $S \seq \RR$ se numește \emph{densă în $\RR$} dacă între oricare
  două numere reale există un element al lui $S$.

  Echivalent, $\forall a < b \in \RR, S \cap (a, b) \neq \emptyset$.
\end{definition}

Rezultatul principal este următorul:

\begin{theorem}\label{thm:q-densa}
  $\QQ$ este densă în $\RR$. Cu alte cuvinte, între oricare două numere reale
  există un număr rațional.
\end{theorem}

\begin{proof}
  Fie $a < b$ două numere reale. Atunci $b - a > 0$, iar din principiul lui Arhimede,
  putem alege $n \in \ZZ$, astfel încît:
  \[
    0 < \frac{1}{n} < b - a \Leftrightarrow a < a + \frac{1}{b} < b.
  \]
  Demonstrația se încheie dacă arătăm că există $m \in \ZZ$, cu $a < \frac{m}{n} < b$.

  Definim:
  \[
    S = \{x \in \ZZ \mid \frac{x}{n} > a \} = \{x \in \ZZ \mid x > na, n > 0 \}.
  \]
  Această mulțime este mărginită inferior de $na$ și este nevidă, deoarece $\ZZ$ nu
  este mărginită superior. Așadar, $S$ are un cel mai mic element, $m$. Arătăm că $m$
  este întregul căutat.

  Cum $\frac{m}{n} > a$, dar $\frac{m - 1}{n} \leq a$, putem scrie:
  \[
    \frac{m}{n} = \frac{m - 1}{n} + \frac{1}{n} \Rightarrow %
    a < \frac{m}{n} = \frac{m - 1}{n} + \frac{1}{n}.
  \]

  Totodată, avem și $a < \frac{m}{n} \leq a + \frac{1}{n} < b$, ceea ce încheie demonstrația.
\end{proof}

%%%%%%%%%%%%%%%%%%%%%%%%%%%%%%%%%%%%%%%%%%%%%%%%%%%%%%%%%%%%%%%%%%%%%%

\section{Mulțimi complete}
\label{sec:complete}

\begin{definition}\label{def:completa}
  O submulțime $S \seq \RR$ se numește \emph{completă} dacă satisface proprietatea supremumului
  (implicit, și pe cea a infimumului).

  Cu alte cuvinte, $S$ este completă dacă orice submulțime nevidă și mărginită a lui $S$
  are și un supremum, și un infimum \emph{în $S$}.
\end{definition}

De exemplu, $[0, 1]$ este completă, $(0, 1)$ nu este completă, iar $\RR$ este completă.

\begin{proposition}\label{pr:q-completa}
  $\QQ$ nu este completă.
\end{proposition}

\begin{proof}
  Fie $S = \{ x \in \QQ \mid 0 < x < \sqrt{2} \} \seq \QQ$.

  Arătăm că $\QQ$ nu este completă, deoarece $S$ nu satisface proprietatea supremumului.

  Cum $\sqrt{2}$ este o margine superioară pentru $S$, avem $\sup S \leq \sqrt{2}$. Presupunem
  inegalitatea strictă. Din densitatea lui $\QQ$ în $\RR$, găsim $q \in \QQ$ cu $\sup S < q < \sqrt{2}$.
  Atunci $q \in S$ și rezultă $q \leq \sup S$, iar nu $\sup S < q < \sqrt{2}$.

  Așadar, $\sup S = \sqrt{2} \notin \QQ$, iar demonstrația este încheiată.
\end{proof}

\end{document}